% Template for PLoS
% Version 3.4 January 2017
\documentclass[10pt,letterpaper]{article}
\usepackage[top=0.85in,left=2.75in,footskip=0.75in]{geometry}

% amsmath and amssymb packages, useful for mathematical formulas and symbols
\usepackage{amsmath,amssymb}

% Use adjustwidth environment to exceed column width (see example table in text)
\usepackage{changepage}

% Use Unicode characters when possible
\usepackage[utf8x]{inputenc}

% textcomp package and marvosym package for additional characters
\usepackage{textcomp,marvosym}

% cite package, to clean up citations in the main text. Do not remove.
% \usepackage{cite}

% Use nameref to cite supporting information files (see Supporting Information section for more info)
\usepackage{nameref,hyperref}

% line numbers
\usepackage[right]{lineno}

% ligatures disabled
\usepackage{microtype}
\DisableLigatures[f]{encoding = *, family = * }

% color can be used to apply background shading to table cells only
\usepackage[table]{xcolor}

% array package and thick rules for tables
\usepackage{array}

% create "+" rule type for thick vertical lines
\newcolumntype{+}{!{\vrule width 2pt}}

% create \thickcline for thick horizontal lines of variable length
\newlength\savedwidth
\newcommand\thickcline[1]{%
  \noalign{\global\savedwidth\arrayrulewidth\global\arrayrulewidth 2pt}%
  \cline{#1}%
  \noalign{\vskip\arrayrulewidth}%
  \noalign{\global\arrayrulewidth\savedwidth}%
}

% \thickhline command for thick horizontal lines that span the table
\newcommand\thickhline{\noalign{\global\savedwidth\arrayrulewidth\global\arrayrulewidth 2pt}%
\hline
\noalign{\global\arrayrulewidth\savedwidth}}


% Remove comment for double spacing
%\usepackage{setspace}
%\doublespacing

% Text layout
\raggedright
\setlength{\parindent}{0.5cm}
\textwidth 5.25in
\textheight 8.75in

% Bold the 'Figure #' in the caption and separate it from the title/caption with a period
% Captions will be left justified
\usepackage[aboveskip=1pt,labelfont=bf,labelsep=period,justification=raggedright,singlelinecheck=off]{caption}
\renewcommand{\figurename}{Fig}

% Use the PLoS provided BiBTeX style
% \bibliographystyle{plos2015}

% Remove brackets from numbering in List of References
\makeatletter
\renewcommand{\@biblabel}[1]{\quad#1.}
\makeatother

% Leave date blank
\date{}

% Header and Footer with logo
\usepackage{lastpage,fancyhdr,graphicx}
\usepackage{epstopdf}
\pagestyle{myheadings}
\pagestyle{fancy}
\fancyhf{}
\setlength{\headheight}{27.023pt}
\lhead{\includegraphics[width=2.0in]{PLOS-submission.eps}}
\rfoot{\thepage/\pageref{LastPage}}
\renewcommand{\footrule}{\hrule height 2pt \vspace{2mm}}
\fancyheadoffset[L]{2.25in}
\fancyfootoffset[L]{2.25in}
\lfoot{\sf PLOS}

%% Include all macros below
\newcommand{\lorem}{{\bf LOREM}}
\newcommand{\ipsum}{{\bf IPSUM}}





\usepackage{forarray}
\usepackage{xstring}
\newcommand{\getIndex}[2]{
  \ForEach{,}{\IfEq{#1}{\thislevelitem}{\number\thislevelcount\ExitForEach}{}}{#2}
}

\setcounter{secnumdepth}{0}

\newcommand{\getAff}[1]{
  \getIndex{#1}{York University}
}

\providecommand{\tightlist}{%
  \setlength{\itemsep}{0pt}\setlength{\parskip}{0pt}}

\begin{document}
\vspace*{0.2in}

% Title must be 250 characters or less.
\begin{flushleft}
{\Large
\textbf\newline{Simple principles for engineering reproducible solutions to
environmental management challenges.} % Please use "sentence case" for title and headings (capitalize only the first word in a title (or heading), the first word in a subtitle (or subheading), and any proper nouns).
}
\newline
\\
Malory Owen\textsuperscript{\getAff{York University}},
Christopher J. Lortie\textsuperscript{\getAff{Biology, York University}}\textsuperscript{*}\\
\bigskip
\textbf{\getAff{York University}}Biology, 4700 Keele St.~Toronto, ON, Canada, M3J1P3\\
\bigskip
* Corresponding author: lortie@yorku.ca\\
\end{flushleft}
% Please keep the abstract below 300 words
\section*{Abstract}
An environmental management challenge is an opportunity to use
fundamental science to inform evidence-based decisions for environmental
stakeholders and conservationists. Contemporary science is embracing
open science and increasingly conscious of reproduciblility.
Synergistically, applying these two paradigms in concert advances our
capacity to move beyond context dependency and singlular thinking to
reverse engineer solutions from published scientific evidence associated
with one challenge to many. Herein, we provide a short list of
principles that can guide those that seek solutions to address
environmental management through primary scientific literature.

% Please keep the Author Summary between 150 and 200 words
% Use first person. PLOS ONE authors please skip this step.
% Author Summary not valid for PLOS ONE submissions.
\section*{Author summary}
Grand challenges require grand solutions. Environmental management
cannot neglect fundamental science as a substrate for effective decision
making, and scientists should be conscious of how their science can be
used by managers.

\linenumbers

% Use "Eq" instead of "Equation" for equation citations.
\section{Introduction}\label{introduction}

Conservation decisions can reside with legislators (indirect impactors)
or environmental managers (direct impactors). To focus on the latter,
managers typically have scientific backgrounds and routinely navigate
the technical literature. However, engagement with scientific literature
is non-trivial for all practitioners because of time, access, relevance
of the science, and reporting standards. Environmental managers and
conservationist are certainly seeking evidence to inform decisions
(citations please). However, there can be a gap in communication between
basic science and management for at least three reasons. Firstly, the
research is not a direct study of an ecosystem, and an immediate,
real-world solution is needed by managers - preferably with a
demonstrable outcome and reasonable cost estimate (citations). This is a
very real limitation in the primary science literature restoration
ecology for instance (cite Lortie better late than never). Secondly, the
link between the biology or ecology present in the literature is not
articulately connected to the similar process for the system at hand.
There are notable examples with journals just as the Journal of Applied
Ecology, Basic and Applied Ecology, the Journal of Environmental
Engineering, and many others, but there nonetheless remains a vast pool
of opportunity for solution development in other journals.
Context-specific findings in science are a legimitate and useful means
to advance discovery, but at times, studies from one system can be
repurposed for insights into another. Finally, the capacity to see the
forest for the trees or the generality of a specific, even large-scale
or broad basic research study can be a challenge. Science can be very
specialized (citation to Nancy Baron book), and mobilizing knowledge for
solutions requires both detailed expertise, scientific synthesis tools
(cite to Lortie formalized synthesis paper in Oikos), or a focus on
identifying the salient elements associated with a study (citation).
Often, seeing the forest also requires sampling many trees. This leads
to the general proposal here that experts can promote reproducible and
mobile knowledge for many fundamental science endeavors by considering
these principles. This is both a set of principles for how to make your
research potentially reusable by environmental managers and
conservationists (citation to get your paper used by meta) and inform
solutions for the environmental crisis.

An environmental management challenge is a problem presented in
scientific literature or society that, when redefined and reviewed using
these 10 principles, can result in a solution to the original problem.
EXPAND\ldots{} a tiny bit more..how. Typically, a challenge is\ldots{}.
look up how they define grand challenges and make mention here including
the 17 global sustainability goals. A (reverse-engineered) reproducible
solution is a suggested solution to a challenge derived from identifying
all the components of the challenge. It should be applicable to multiple
local-extent challenges when tweaked to fit the circumstances. Finally,
a tool or solution is the desired outcome from the primary research to
support evidence-based decision making in conservation. In this case, a
tool is a methodology researchers use that can faciitate managers to
either identify best ways to measure/identify issues or to provide
solutions for their specific challenge. Any tool is linked to its
respective reproducible solution by the fundamental concepts of
reproducibility (citation to reproducibility crisis paper). These can
include primarily conceptual replication but also\ldots{}. check papers
on reproducibility. then concluding sentence.

The heuristic developed here was inspired by the ``Ten Simple Rules''
format pioneered by Phillip Bourne in the field of computational biology
{[}1{]}. We propose that by distilling the concepts that promote
engagement with scientific literature outside of the research community,
managers can rely on broader sources of scientific knowledge to make
decisions. Additionally, researchers can better understand the
perspective of managers facilitating science and scientific
communication that is more applicable to managers without compromising
their research programs. Here, we will outline and discuss simple
``principles'' scientists can use to make their research more applicable
to managers, and managers can use to identify basic science that fits
their needs.

\section{The principles}\label{the-principles}

\textbf{1. Reframe the problem as challenge.} Doom-and-gloom is a
pervasive mentality in ecology and environmental sciences which can
reduce productivity and problem solving. It can shut down even the most
motivated of minds--but beyond the issue of motivation, reframing a
problem as a \emph{challenge} can reveal solution sets that may
otherwise remain hidden. For example, consider the problem of
human-wildlife conflict between carnivores and the people living near
the Ruaha National Park boundary in Tanzania. The \emph{problem} is that
98.5\% of people perceive wildlife as a threat to their livestock
resulting in increased likelihoods for human-wildlife conflict {[}2{]}.
Reframed, the challenge can be to decrease the negative perception of
wildlife in areas with high human-wildlife conflict opportunities. It is
a small change in semantics, but it is a potentially profound change in
direction.

\textbf{2. Describe the scope and extent of the challenge.} Defining the
scope of a challenge conceptually and the extent geographically will
ensure that potential solutions fit the challenge. Moving across scales
is a common issue in ecology (citation to Sandel 2015 paper in
Ecography), and proposing a spatial scale, using common terms, and
describing the breadth of the challenge will accelerate
interdisciplinary solutions (i.e.~the wildlife-human challenge above is
ecological and societal). The challenge can be problematic on multiple
scales including local, regional, or global scale and solutions can be
needed for each. Conceptually, the scope is broad in the human-wildlife
conflict example whilst the extent is primarily local to the area
surrounding the Southern border of the Ruaha National Park.

\textbf{3. Explicitly link the basic science to management implications
and policy.} Perhaps the most facile principle, a simple description and
definition of the basic scientific evidence in a study and how it can be
resued is a fundamental step in linking science to evidence-based
decision making for environmental challenges. In the wildlife-human
challenge, depredation of livestock impacted 61.1\% of households, but
livestock loss due to disease or theft was X percent? of total loss
{[}2{]}. Consequently, a clear statement of evidence can illuminate the
most viable solution sets in some instances.

\textbf{4. Propose implications of ignoring this challenge.} A
description of the impact a challenge on a system if left unchecked will
help clarify the severity of the challenge. The trickle down effects and
indirect implications of the challenge should also be examined. For
instance, anti-carnivore sentiment will likely only grow as climate
change and pressures to confine pastoral herders makes livestock more
difficult to raise {[}3,4{]}. Many large carnivores are already
threatened/endangered, and further anthropogenic pressures on the
populations will lead to severe declines in populations including
potential extinction of keystone species {[}5--7{]}; but it is often
associated with underlying human-human conflict {[}8{]}.

\textbf{5. State the direct human needs associated with this challenge.}
To expand on principle 4, considering the direct needs of humans is an
essential part of the process. Conservationists have a tendency to react
sentimentally, considering the intrinsic value of the ecosystem and its
right to exist. This is valid, and should not be overlooked. However, it
is all too easy for those who value ``nature'' to dismiss ecosystem
services. Identifying anthropogenic needs will help a problem solver
create a solution that is appropriate for the challenge, and prevent the
manifestation of new challenges. To bring it back to our human-wildlife
conflict challenge, a direct need of humans in the Ruaha region is to
have livestock with which to make a living. It is true that a solution
\emph{could} be to remove livestock; that way, the carnivores will no
longer be seen as predators of non-existent livestock. Obviously, this
is not a viable solution as the humans interacting within this system
will suffer, rendering the solution unethical and unsustainable. The
error in this example may seem obvious, but explicitly identifying all
needs associated with a challenge will help organize the problem-solving
process, and prevent any potential missteps. It is always a good idea to
ask the group affected--particularly if they are a marginalized
group--what their direct needs are as well.

\textbf{6. List at least one limitation of the study and explain.} While
the peer review process should eliminate any poorly executed science,
nearly all studies have some limitation. Critically reading the study
associated with the challenge can mean the difference between success
and failure of a later implemented management solution that otherwise
follows all other principles presented. For example, conflating
causation with correlation can be a fatal misstep if making real-world
decisions based on a correlated trend. However, it's an unfortunate
truth that conservationists and managers often have to make decisions on
a real-world challenge without fully understanding the system in which
they work, since a solution is often needed immediately in order to
prevent a catastrophic loss. Still, knowing the relative strength of a
solution can help prepare for a range of outcomes. Perhaps the most
important step, however, is to identify limitations of the study in
order to best advocate for an eventual solution's implementation. Being
well versed in counter-arguments is essential in any debate, and it is
no different when advocating for a solution.

\textbf{7. Explore the benefits of minimal intervention for
stakeholders.} The culmination of principles of 4, 5, and 6, principle 7
is another idea that can feel uncomfortable to a traditional
conservationist. Why not just do \emph{everything} so that all our bases
are covered? Unfortunately, resources to conserve the environment are
limited, and managers must maximize their ``profits'' (positive
restoration/conservation outcomes). Ideally, a well-designed solution
takes minimal resource allocation, and has minimal incidental negative
effects.

\textbf{8. List the tools applied to this challenge.} This principle is
most critical for researchers wishing to make their science more
amenable to management use, but is also helpful to managers seeking
relevant papers beyond studies which focus specifically on their
projects. In an environmental management challenge case study, there is
usually one primary tool the researchers use to produce their results.
However, there can be many more than one! These tools usually take the
form of common methodologies that can be used across sub-disciplines,
and are easily identifiable. Look to the methods to identify these
tools. Some examples of primary tools we have found include:
\emph{meta-analysis {[}9{]}; {[}10{]} }systematic review {[}11{]}
\emph{citizen science {[}12{]}; {[}13{]} }team science {[}14{]} \emph{R
{[}15{]} }mapping {[}16{]} \emph{big data/open access databases
{[}17{]}; {[}18{]}; {[}19{]} }modelling {[}20{]} \emph{surveys {[}21{]}
}biomonitoring {[}22{]} \emph{evolutionary change/population viability
genetic analysis {[}23{]} }economic incentivization monitoring {[}24{]}
*human health and well-being monitoring {[}25{]}; {[}2{]} This list
could potentially be expanded upon, and new tools can result in new
solutions.

\textbf{9. Explain the role that the primary tool addressed for the
challenge--i.e.~identification/research evidence, management/solution
applied, or inform policy.} Continuing from principle 8, this principle
focuses on which tool used in a study is primarily responsible for
producing results. A good place to start may be the tile, as often a
tool will be highlighted; for example, ``Odonata (Insecta) as a tool for
the bio-monitoring of environmental quality'' {[}22{]}. After
identifying which tool is primary, consider the role that tool could
play in finding the solution to the challenge. Three examples of roles a
tool could play are: \emph{identification/research evidence} which seeks
to better our understanding of systems so that we may more effectively
manage them. It is the most ``pure'' role of tools in basic science, and
is also the most common in typical ecology papers. Tools that could play
this role include surveys, meta analysis, team science, systematic
review, citizen science, R, mapping, open access databases, modelling.
\emph{management/solution applied} which observes the effectiveness of
management strategies. Tools that could play this role include
bio-monitoring and evolutionary change/population viability genetic
analysis. \emph{inform policy} which seeks to help legislators make
decisions about environmental challenges. They often take the form of
anthropocentric studies that consider ecological or environmental
policy. Any tool could play this role, but some tools which fit most
neatly into this role include economic incentivization monitoring and
human health and well-being monitoring. Depending on how a tool is used,
these tools may play multiple roles in a solution, or may play an
entirely different role than described. Most importantly, the primary
tool's role informs the direction our solution should take, based on our
challenge study.

\textbf{10. Apply the tool to another challenge or explain how it is
general and scaleable.} Finally, use the primary tool explored in
principle 9 to broaden the applicability of the challenge's solution.
This principle ties together the idea that reverse-engineered
reproducible solutions are relevant to more than just the environmental
management challenge case unpacked. After doing all this work, it makes
sense to seek and understand its reproducibility for future work! Beyond
efficiency, it also addresses the key challenge we are trying to address
through these principles: Basic science can be more applicable than it
sometimes appears.

\section{Implications}\label{implications}

At its core, our 10 principles for reverse-engineering solutions to
solve environmental challenges are guidelines to improve science
communication. Usually, when we talk about sci comm, the burden of
effective communication is on the researcher or advocate (because
usually we are directing our sci comm at general audiences who are not
obligated to give the researcher their attention). But in animal
behavior, communication requires at least two parties: a \emph{sender}
and a \emph{reciever}. In this instance, the researcher is a sender and
the manager is the receiver. For this communication to be successful,
both parties must actively participate in this reuse of solutions from
one domain to another. Luckily, both parties in our scenario have
incentive to engage with effective communication. Professional advocacy
is a common priority for ecologists, who hope their work influences
better decision-making of both indirect and direct impactors. And
managers, of course, wish to improve their project outcomes. Increased
consumption of scientific evidence by managers, and basic science that
is more palatable to a broader audience written by researchers, will
result in \emph{functional use of scientific literature}. Overall,
better collaboration will encourage the realization of the common goal:
solving an environmental management challenge.

\section*{References}\label{references}
\addcontentsline{toc}{section}{References}

\hypertarget{refs}{}
\hypertarget{ref-Bourne2006}{}
1. Bourne PE, Chalupa LM. Ten simple rules for getting grants. PLoS
Computational Biology. 2006;2: 59--60.
doi:\href{https://doi.org/10.1371/journal.pcbi.0020012}{10.1371/journal.pcbi.0020012}

\hypertarget{ref-Dickman2014}{}
2. Dickman AJ, Hazzah L, Carbone C, Durant SM. Carnivores, culture and
'contagious conflict': Multiple factors influence perceived problems
with carnivores in Tanzania's Ruaha landscape. Biological Conservation.
Elsevier Ltd; 2014;178: 19--27.
doi:\href{https://doi.org/10.1016/j.biocon.2014.07.011}{10.1016/j.biocon.2014.07.011}

\hypertarget{ref-Jones2009}{}
3. Jones PG, Thornton PK. Croppers to livestock keepers: livelihood
transitions to 2050 in Africa due to climate change. Environmental
Science and Policy. 2009;12: 427--437.
doi:\href{https://doi.org/10.1016/j.envsci.2008.08.006}{10.1016/j.envsci.2008.08.006}

\hypertarget{ref-Lindsey2009}{}
4. Lindsey PA, Romañach SS, Davies-Mostert HT. The importance of
conservancies for enhancing the value of game ranch land for large
mammal conservation in southern Africa. Journal of Zoology. 2009;277:
99--105.
doi:\href{https://doi.org/10.1111/j.1469-7998.2008.00529.x}{10.1111/j.1469-7998.2008.00529.x}

\hypertarget{ref-Towns2009}{}
5. Towns L, Derocher AE, Stirling I, Lunn NJ, Hedman D. Spatial and
temporal patterns of problem polar bears in Churchill, Manitoba. Polar
Biology. 2009;32: 1529--1537.
doi:\href{https://doi.org/10.1007/s00300-009-0653-y}{10.1007/s00300-009-0653-y}

\hypertarget{ref-Bagchi2006}{}
6. Bagchi S, Mishra C. Living with large carnivores: Predation on
livestock by the snow leopard (Uncia uncia). Journal of Zoology.
2006;268: 217--224.
doi:\href{https://doi.org/10.1111/j.1469-7998.2005.00030.x}{10.1111/j.1469-7998.2005.00030.x}

\hypertarget{ref-Johnson2006}{}
7. Johnson A, Vongkhamheng C, Hedemark M, Saithongdam T. Effects of
human-carnivore conflict on tiger (Panthera tigris) and prey populations
in Lao PDR. Animal Conservation. 2006;9: 421--430.
doi:\href{https://doi.org/10.1111/j.1469-1795.2006.00049.x}{10.1111/j.1469-1795.2006.00049.x}

\hypertarget{ref-Dickman2010}{}
8. Dickman AJ. Complexities of conflict: The importance of considering
social factors for effectively resolving human-wildlife conflict. Animal
Conservation. 2010;13: 458--466.
doi:\href{https://doi.org/10.1111/j.1469-1795.2010.00368.x}{10.1111/j.1469-1795.2010.00368.x}

\hypertarget{ref-Castanho2015}{}
9. Castanho C de T, Lortie CJ, Zaitchik B, Prado PI. A meta-analysis of
plant facilitation in coastal dune systems: responses, regions, and
research gaps. PeerJ. 2015;3: e768.
doi:\href{https://doi.org/10.7717/peerj.768}{10.7717/peerj.768}

\hypertarget{ref-Busch2017}{}
10. Busch J, Ferretti-Gallon K. What drives deforestation and what stops
it? A meta-analysis. Review of Environmental Economics and Policy.
2017;11: 3--23.
doi:\href{https://doi.org/10.1093/reep/rew013}{10.1093/reep/rew013}

\hypertarget{ref-Lortie2018}{}
11. Lortie CJ, Filazzola A, Kelsey R, Hart AK, Butterfield HS. Better
late than never: a synthesis of strategic land retirement and
restoration in California. Ecosphere. 2018;9: e02367.
doi:\href{https://doi.org/10.1002/ecs2.2367}{10.1002/ecs2.2367}

\hypertarget{ref-Burkle2013}{}
12. Burkle LA, Marlin JC, Knight TM. Plant-Pollinator Interactions over
120 Years: Loss of Spcies, Co-Occurence, and Function. Science.
2013;339: 1611--1616.

\hypertarget{ref-Conrad2011}{}
13. Conrad CC, Hilchey KG. A review of citizen science and
community-based environmental monitoring: Issues and opportunities.
Environmental Monitoring and Assessment. 2011;176: 273--291.
doi:\href{https://doi.org/10.1007/s10661-010-1582-5}{10.1007/s10661-010-1582-5}

\hypertarget{ref-Nielsen2017}{}
14. Nielsen JA, Grøndahl E, Callaway RM, Dickinson KJ, Ehlers BK. Home
and away: biogeographical comparison of species diversity in Thymus
vulgaris communities. Biological Invasions. Springer International
Publishing; 2017;19: 2533--2542.
doi:\href{https://doi.org/10.1007/s10530-017-1461-x}{10.1007/s10530-017-1461-x}

\hypertarget{ref-McCarthy2010}{}
15. McCarthy MP, Best MJ, Betts RA. Climate change in cities due to
global warming and urban effects. Geophysical Research Letters. 2010;37:
1--5.
doi:\href{https://doi.org/10.1029/2010GL042845}{10.1029/2010GL042845}

\hypertarget{ref-Halpern2008}{}
16. Halpern B, Walbridge S, Selkoe K, Kappel C, Micheli F, D'Argrosa C,
et al. A Global Map of Human Impact on Marine Ecosystems. Science.
2008;319: 948--953.
doi:\href{https://doi.org/10.1111/2041-210X.12109}{10.1111/2041-210X.12109}

\hypertarget{ref-Sillero2014}{}
17. Sillero N, Campos J, Bonardi A, Corti C, Creemers R, Crochet PA, et
al. Updated distribution and biogeography of amphibians and reptiles of
Europe. Amphibia Reptilia. 2014;35: 1--31.
doi:\href{https://doi.org/10.1163/15685381-00002935}{10.1163/15685381-00002935}

\hypertarget{ref-Dengler2011}{}
18. Dengler J, Jansen F, Glöckler F, Peet RK, Cáceres M de, Chytrý M, et
al. The Global Index of Vegetation-Plot Databases (GIVD): A new resource
for vegetation science. Journal of Vegetation Science. 2011;22:
582--597.
doi:\href{https://doi.org/10.1111/j.1654-1103.2011.01265.x}{10.1111/j.1654-1103.2011.01265.x}

\hypertarget{ref-Maldonado2015}{}
19. Maldonado C, Molina CI, Zizka A, Persson C, Taylor CM, Albán J, et
al. Estimating species diversity and distribution in the era of Big
Data: To what extent can we trust public databases? Global Ecology and
Biogeography. 2015;24: 973--984.
doi:\href{https://doi.org/10.1111/geb.12326}{10.1111/geb.12326}

\hypertarget{ref-Vogt2018}{}
20. Vogt R, Sharma S, Leavitt P. Direct and interactive effects of
climate, meteorology, river hydrology, and lake characteristics on water
quality in productive lakes of the Canadian Prairies. Canadian Journal
of Fisheries and Aquatic Sciences. 2018;75: 47--59.
doi:\href{https://doi.org/10.1139/cjfas-2016-0520}{10.1139/cjfas-2016-0520}

\hypertarget{ref-Wassen2005}{}
21. Wassen MJ, Venterink HO, Lapshina ED, Tanneberger F. Endangered
plants persist under phosphorus limitation. Nature. 2005;437: 547--550.
doi:\href{https://doi.org/10.1038/nature03950}{10.1038/nature03950}

\hypertarget{ref-Miguel2017}{}
22. Miguel TB, Oliveira-Junior JMB, Ligeiro R, Juen L. Odonata (Insecta)
as a tool for the biomonitoring of environmental quality. Ecological
Indicators. Elsevier; 2017;81: 555--566.
doi:\href{https://doi.org/10.1016/j.ecolind.2017.06.010}{10.1016/j.ecolind.2017.06.010}

\hypertarget{ref-Stoops2016}{}
23. Stoops MA, Campbell MK, DeChant CJ, Hauser J, Kottwitz J, Pairan RD,
et al. Enhancing captive Indian rhinoceros genetics via artificial
insemination of cryopreserved sperm. Animal Reproduction Science.
Elsevier B.V. 2016;172: 60--75.
doi:\href{https://doi.org/10.1016/j.anireprosci.2016.07.003}{10.1016/j.anireprosci.2016.07.003}

\hypertarget{ref-Cerda2018}{}
24. Cerda C, Fuentes JP, De La Maza CL, Louit C, Araos A. Assessing
visitors' preferences for ecosystem features in a desert biodiversity
hotspot. Environmental Conservation. 2018;45: 75--82.
doi:\href{https://doi.org/10.1017/S0376892917000200}{10.1017/S0376892917000200}

\hypertarget{ref-Mergler2007}{}
25. Mergler D, Anderson HA, Chan LHM, Mahaffey KR, Murray M, Sakamoto M,
et al. Methylmercury exposure and health effects in humans: A worldwide
concern. Ambio. 2007;36: 3--11.
doi:\href{https://doi.org/10.1579/0044-7447(2007)36\%5B3:MEAHEI\%5D2.0.CO;2}{10.1579/0044-7447(2007)36{[}3:MEAHEI{]}2.0.CO;2}

\nolinenumbers


\end{document}

